\section{General complex exclusion of \texorpdfstring{$\sha{R8_3_3472}$}{R8-3-3472}}
\label{sec:r3472-focus}
The remaining case $R=\sha{R8_3_3472}$ is excluded by a complete
central-induction and rotation argument. The earlier feasible center
relaxations and conditional Morita restrictions are retained below with
their historical scope. They are not the terminal obstruction.

\begin{theorem}\label{thm:r3472-general-exclusion}
The fusion ring $\sha{R8_3_3472}$ has no fusion categorification over
$\mathbb C$. In particular, the conclusion does not require a pivotal,
spherical, braided, unitary, or pseudo-unitary structure as an initial
hypothesis.
\end{theorem}
\begin{proof}
Theorem~\ref{thm:research-four-general-spherical} reduces a hypothetical
complex categorification to a pseudo-unitary Galois conjugate with the
same fusion rules. We recall that reduction, then exclude every possible
center restriction frame. There are $372$ labelled low frames; exact
span and phase certificates leave three equivalent frames. Their twists
belong to a complete domain of $424$ Galois orbit types. The necessary
rotation equations for these types have an exact integer contradiction,
proved by $1648$ rational bound deductions and a final separating
inequality. The details and replay scope follow.
\end{proof}

\subsection{Reduction from arbitrary complex categorification}
\label{subsec:r3472-exclusion-scope}
All eight source simples are selfdual. In the canonical sphericalization
of an arbitrary categorification, both lifts of each simple are selfdual,
and changing the lift reverses its second indicator
\cite[Section~5.1]{OstrikPiv3}. Choose the lift of indicator $+1$.
The odd-multiplicity theorem~\cite[Theorem~3.9]{LPRV} makes the
multiplicity of every negative lift in a positive-lift square even.
Consequently the signed quotient tensor satisfies
\[
 |T_{ij}^k|\leq N_{ij}^k,\qquad
 T_{ij}^k\equiv N_{ij}^k\pmod2,\qquad
 T_{jj}^i\equiv N_{jj}^i\pmod4.
\]
The complete search, with the unit, duality, cyclic trace symmetry,
dual reversal and all ordered associativity equations, has $84$ signed
variables and $298865973$ nodes. Its only solution is $T=N$.
No further sign gauge or commutativity assumption is added. The positive
lifts therefore form a spherical section, proving that the original
category admits a spherical structure.

Of the eight ordinary characters of $R$, the six outside the FP Galois
orbit have zero values on simple basis elements. They cannot be spherical
dimension characters. The remaining pair is the FP character and its
$\sqrt2$-conjugate. Extend the required embedding to an automorphism of
$\mathbb C$ and transport the associator, rigidity and spherical structure
simultaneously. All coherence identities are preserved, and the resulting
spherical dimensions are positive FP dimensions. Thus the transported
category is pseudo-unitary. This also avoids any separate assumption that
a chosen pivotal structure descends to a number field.
The original gate, source hashes and complete replay remain in
\path{evidence/rank8_research_20261022/pivotality/}; the new proof package
contains an independent audit and fresh successful replay of this step.

\subsection{The complete distinguished restriction frames}
Put
\[
 r=\sqrt2,\quad \delta=7+5r,\quad
 \pi=\delta^2=99+70r,\quad \beta=1+\pi=100+70r,
 \quad \gamma=7\delta=49+35r,\quad D=7\beta.
\]
The source dimensions, in canonical order, are
$(1,\delta-1,\delta-1,\delta,\delta-1,\delta+1,\delta+1,\delta+1)$.
Let $V=X_1+\cdots+X_7$. We use $F:Z(\mathcal C)\to\mathcal C$ and
its adjoint $I$, and write $v_Z$ for the nonnegative integral restriction
column of a central simple $Z$.
The source formal codegrees are $D,D'=700-490r$, and six copies of $7$.
Central induction gives eight distinct twist-one summands $A_E$ of
$I(\one)$, with dimensions $1,\pi,\beta,\ldots,\beta$
\cite[Theorem~2.13]{OstrikPiv3}.

The restriction of the distinguished object of dimension $\pi$ is forced:
\begin{equation}\label{eq:r3472-fixed-pi}
 F(A_{D'})=\one+2V.
\end{equation}
To avoid an identification assumption, first let $P$ be the image of the
vacuum under a modular Galois automorphism that flips $r$. The Galois
vacuum row gives $\dim(P)=D/D'=\pi$ and, by the induction character
formula~\cite[Lemma~4.3 and Corollary~4.5]{DeshpandeCenter},
\[
 \chi_E(F(P))=\begin{cases}
 \pi&E=\mathrm{FP},\\
 \pi^{-1}&E=\mathrm{FP}',\\
 1&f_E=7.
 \end{cases}
\]
The element $\one+2V$ has exactly these eight character values.
Character inversion proves $F(P)=\one+2V$. Its unit multiplicity is one,
so $P$ belongs to $I(\one)$ and is the distinguished $A_{D'}$.
This proves \eqref{eq:r3472-fixed-pi} without presuming that no other
central object has dimension $\pi$.

The induction Gram matrix is the tensor-derived integer matrix
$B=\sum_iN_iN_i^{\mathsf T}=\sum_Zv_Zv_Z^{\mathsf T}$.
Set $u=e_0$, $p=(1,2,\ldots,2)^{\mathsf T}$ and
$Q=B-uu^{\mathsf T}-pp^{\mathsf T}$. A distinguished $\beta$ column
$v$ has $v_0=1$ and $\sum_{i>0}v_i=14$. With
$A=v_1+v_2+v_4$ and $C=v_5+v_6+v_7$, its dimension equation gives
$C=A+1$ and $v_3=13-2A$, where $0\leq A\leq6$.
These equations give exactly $2142$ possible columns. The necessary
condition $Q-vv^{\mathsf T}\succeq0$ retains $96$.
For the six labelled columns, impose distinguished $S$-symmetry, the
pairwise Gram inequalities, their prescribed sum $Q_{0,*}$, and positivity
of the residual Gram matrix. Exact exhaustive enumeration leaves $372$
labelled frames. Repeated restriction columns are allowed throughout.
A separate enumeration and exact positive-semidefinite test reproduce
this complete list.

The high columns, those with $v_0=0$, are reconstructed from the source
rather than taken from a previously feasible witness. The complete
induction bounds, modular eigenvalue integrality and Galois spectral
support tests leave $90$ restrictions of dimension
$\eta=\beta/2$, $78$ of dimension $\beta$, and one each of dimensions
$c\gamma$ for $c=1,2,5,10$; the latter columns are $cV$.
Let $a,b$ count the high $\eta,\beta$ simples, and put
$w=\sum c^2$ over the other high simples. Comparing global dimension
and orthogonality of the dimension character with its $r$-conjugate gives
\[
 50a+200b+49w=8402,\qquad
 50a+200b-49w=-1202.
\]
Hence $a+4b=72$ and $w=98$, excluding $c=10$.
The resulting $171$ high restriction possibilities form a complete
necessary domain. The source-derived domain replay checks each
integrality and missing-Galois-conjugate rejection exactly.

\subsection{Reduction of all frames to one orbit}
Let $X=(\chi_E(X_j))$ be the full character table and let $L$ have
columns $v_{A_E}$ in the same character order. The distinguished
modular $STS$ identities and induction give
\begin{equation}\label{eq:r3472-full-sts-restriction}
 \sum_Z\theta_Zv_Zv_Z^{\mathsf T}=LX.
\end{equation}
For each frame, the high columns must lie in the support of its residual
positive-semidefinite Gram matrix. Exact linear annihilators show that
\eqref{eq:r3472-full-sts-restriction} cannot lie in the span of the
supported rank-one matrices for $276$ frames.
Nine verified based-ring automorphisms partition the remaining $96$
frames into twelve orbits. Eleven exact rational separating certificates
exclude eleven orbits. The only frames left are $42,117,304$, one orbit
under these automorphisms.

These phase certificates do not assume a maximum twist order. Apply
normalized relative trace to
$K=\mathbb Q(r,\zeta_7+\zeta_7^{-1})$ in
\eqref{eq:r3472-full-sts-restriction}. The $33$ exact trace values of
roots whose orders divide $280$, together with their convex hull, contain
the normalized relative trace of every root of unity. Cyclotomic tower
transitivity and the factor $\mu(m)/\varphi(m)$ for the part of the
order coprime to $14$ prove the coverage. Introducing nonnegative real
masses on these atoms is therefore a necessary relaxation for arbitrary
twists. Each separating covector is checked against every supported
column and atom, as well as the exact right-hand side. The frame-coverage
replay verifies all $276$ span certificates, the nine automorphisms,
all eleven phase certificates, and the three surviving labels.

\subsection{Rigidity of the surviving frame}
Use frame $42$, and identify high columns with their seven nonunit
coordinates. Its residual Gram matrix has null-vector
$(1,1,-6,1,1,1,1)^{\mathsf T}$. Thus every high column satisfies
$\sum_{i>0}v_i=7v_3$. An $\eta$ column would give $v_3=1$,
$A+C=6$ and $C-A=1$, impossible over the integers. Hence there are
$18$ high $\beta$ simples and $\sum c^2=98$ for the remaining objects.
For a high $\beta$ column, $v_3=2$, $A=5$, and $C=7$.
Writing $z=v-2V$ gives $\|z\|^2\geq2$. The residual Gram trace is
$1226$ and forces $\sum\|z\|^2=36$, so equality holds for all
$18$ columns. Consequently the only possibilities are
\[
 2V-e_a+e_c,\qquad a\in\{1,2,4\},\quad c\in\{5,6,7\}.
\]
The cross entries of the Gram matrix force each of these nine columns to
occur twice. The remaining columns are $cV$, $c\in\{1,2,5\}$,
with dimensions $c\gamma$.

The ten associated rank-one matrices are linearly independent.
Solving \eqref{eq:r3472-full-sts-restriction} makes the twist sum on
each two-object $\beta$ column equal to $2\cos(2\pi k/7)$, where,
in the displayed row and column orders, the $k$-matrix is
\[
 \begin{pmatrix}1&3&2\\2&1&3\\3&2&1\end{pmatrix}.
\]
A pair of unit complex numbers with this sum is
$\{\zeta_7^k,\zeta_7^{-k}\}$. Thus the $18$ high $\beta$ twists
are fixed, while the six distinguished $\beta$ twists are one.
For the other simples define
\[
 P_n=\sum_{\dim Z=c\gamma}c^2\theta_Z^n.
\]
The same $STS$ computation and the second and third indicators give
\begin{equation}\label{eq:r3472-first-three-moments}
 P_0=98,\qquad P_1=-2-10r,\qquad P_2=-2,\qquad P_3=-2+10r.
\end{equation}
For the second indicators, odd-square witnesses for objects
$1,2,3,4,5,6,7$ are respectively squares of objects
$4,1,1,2,6,7,5$; Theorem~3.9 of~\cite{LPRV} gives $\nu_2=1$.
The third indicators are all two by
Proposition~\ref{prop:r3472-maximal-third-indicator} below.
The central indicator formula~\cite{NgIndicators} now becomes, for
$i>0$,
\begin{equation}\label{eq:r3472-indicator-moments}
 \nu_n(X_i)=b_i(n)+\frac{r}{20}(P_n+2),\qquad
 b_i(n)=\begin{cases}
 1&7\nmid n,\\
 6&i\in\{1,2,4\},\ 7\mid n,\\
 7&i=3,\ 7\mid n,\\
 8&i\in\{5,6,7\},\ 7\mid n.
 \end{cases}
\end{equation}

\subsection{A complete finite twist domain}
Cauchy's theorem~\cite[Theorem~3.9]{BNRWCauchy} and
$DD'=9800$ restrict exponent primes to $2,5,7$.
Modular Galois symmetry sends a twist $\theta$ to $\theta^{a^2}$
for every unit $a$ modulo the exponent. It preserves each dimension
$c\gamma$: under the $r$-flip, $D/\sigma(D)=\pi$ and
$\sigma(\gamma)=-\gamma/\pi$, with the Galois sign correcting the
minus sign. Therefore the weighted multiplicities in $P_n$ are constant
on unit-square orbits.

For an order $q=2^a5^b7^c$, the orbit size is the product of
\[
 s_2(a)=2^{\max(a-3,0)},\qquad
 s_5(b)=\begin{cases}1&b=0,\\2\cdot5^{b-1}&b>0,\end{cases}
 \quad
 s_7(c)=\begin{cases}1&c=0,\\3\cdot7^{c-1}&c>0.\end{cases}
\]
Since the entire weighted mass is $98$, every occurring orbit has size
at most $98$. It follows that $a\leq9$, $b\leq3$, $c\leq2$, with
the product bound still imposed. There are exactly $63$ possible orders
and $424$ unit-square orbit types. This is a proved finite domain, not a
conductor cutoff. For each type $\mathcal O$ introduce its nonnegative
integer weighted multiplicity $x_{\mathcal O}$; a categorical realization
would have $x_{\mathcal O}=\sum_c c^2m_{c,\mathcal O}$. Allowing every
nonnegative integer $x_{\mathcal O}$ only weakens the necessary system.
Then
\[
 P_n=\sum_{\mathcal O}x_{\mathcal O}
             \sum_{\theta\in\mathcal O}\theta^n,
 \qquad \sum_{\mathcal O}|\mathcal O|x_{\mathcal O}=98.
\]
The power sums are evaluated by exact local Gauss periods. They lie in
$\mathbb Q(\zeta_8,\sqrt5,\sqrt{-7})$; computation in
$\mathbb Q(\zeta_{1120})$ is only a common coordinate representation,
not a restriction that twists have order dividing $1120$.
Independent reduction modulo cyclotomic polynomials checks the local
formulas and all $424$ orbit types.

\subsection{Rotation constraints and the exact terminal contradiction}
Let $E_n$ be pivotal cyclic rotation on
$\operatorname{Hom}(\one,X_i^{\otimes n})$. It has order dividing $n$.
For $1\leq k<n$, set $g=\gcd(n,k)$ and $m=n/g$. The rotation identity
\cite[Theorem~5.1 and its proof]{NgHigherPivotal} gives
\[
 \operatorname{Tr}(E_n^k)
   =\sigma_{k/g}\bigl(\nu_m(X_i^{\otimes g})\bigr),
 \qquad \sigma_{k/g}(\zeta_m)=\zeta_m^{k/g}.
\]
The tensor determines $X_i^{\otimes g}$ exactly, and the indicator is
additive. Thus \eqref{eq:r3472-indicator-moments} computes each trace.
The Fourier multiplicities
\[
 \frac1n\sum_{k=0}^{n-1}\zeta_n^{-tk}\operatorname{Tr}(E_n^k)
 \quad(0\leq t<n)
\]
are nonnegative integers. Impose the necessary integrality congruences
from their constant coefficient in the chosen cyclotomic power basis,
for $i=1,3,5$ and
$n\in\{4,5,7,8,10,16,20\}$. Together with
\eqref{eq:r3472-first-three-moments}, this gives $362$ integer equations
in the $424$ orbit variables and $148$ integer quotient variables.
All variable bounds follow from the mass $98$ and these congruences.

There is a further exact necessary condition:
$\nu_n(X_i)\in\mathbb Q(\zeta_n)$.
If $a$ is coprime to the exponent, $a\equiv1\pmod n$, and
$\chi(a)$ is the sign of the action on $r$,
\eqref{eq:r3472-indicator-moments} implies
\[
 P_n-\chi(a)P_{na}=2(\chi(a)-1).
\]
Their coordinate equations add $266$ rows, giving a necessary system
\begin{equation}\label{eq:r3472-integer-system}
 A x=b,\qquad \ell\leq x\leq u,\qquad x\in\mathbb Z^{572},
 \qquad A\in\mathbb Z^{628\times572}.
\end{equation}
Every coefficient was independently reconstructed with arbitrary-precision
arithmetic from the tensor, exact power sums and rotation identities.

The terminal proof uses only rational arithmetic and integer rounding.
For any rational row covector $y$ and coordinate vector $e_j$, the
identity $Ax=b$ gives the valid lower bound
\[
 x_j\geq yb+
 \sum_k\min\bigl((e_j-yA)_k\ell_k,(e_j-yA)_ku_k\bigr).
\]
Rounding its right-hand side upward is valid because $x_j$ is an integer;
applying the same argument to $-x_j$ gives upper bounds.
The certificate records $380+486+782=1648$ such sequential deductions.
For the resulting box $[\ell^*,u^*]$, its final rational covector $y$ satisfies
\[
 yb-\sum_j\max\bigl((yA)_j\ell_j^*,(yA)_ju_j^*\bigr)
 =\frac{2717541156573712541191977062303}
        {359103652930042901916159999600}>0.
\]
This contradicts $Ax=b$. It excludes frame $42$ and its two equivalents,
so no low frame can be the center of a pseudo-unitary categorification.
The general complex reduction proves
Theorem~\ref{thm:r3472-general-exclusion}.

\subsection{Proof record, replay and historical scope}
The canonical proof package is
\begin{quote}\small\path{evidence/r3472_exclusion_20261022/}.\end{quote}
It binds the canonical source tensor, complete domain and frame audits,
source-derived rotation system, rational cuts and final separator.
The complete replay command is
\begin{quote}\small
\texttt{python3 }\path{evidence/r3472_exclusion_20261022/replay.py}.
\end{quote}
For the self-contained terminal arithmetic certificate alone, run
\begin{quote}\small
\texttt{python3 }\path{evidence/r3472_exclusion_20261022/terminal/check.py}.
\end{quote}
The terminal checker uses Python's standard library and exact fractions;
no optimizer or floating-point feasibility verdict is used as proof.
An independently implemented checker replays the same contradiction.
A separate exact SCIP certificate is also accepted by the unmodified
VIPR checker, with its input independently matched to the audited system.
These arithmetic checks certify the supplied system; the source derivation,
complete frame coverage and general sphericalization reduction are distinct
parts of the full proof and are all included in the record.

The exact-M3 totals are therefore $3507$ general complex exclusions,
$16$ unitary models and four unresolved rings:
$\sha{R8_3_928}$, $\sha{R8_3_941}$, $\sha{R8_3_1154}$ and
$\sha{R8_3_1768}$. Cumulatively through multiplicity three the totals are
$4403$ exclusions, $93$ models and four unresolved rings.
The earlier structural arguments below remain useful records of the
investigation. Their feasible necessary-condition witnesses and failed
construction routes were not nonexistence proofs; the exclusion above is
independent of the existence of either proposed rank-two module category.

\section{Historical Morita and construction restrictions for \texorpdfstring{$\sha{R8_3_3472}$}{R8-3-3472}}
\label{sec:r3472-historical-focus}
This earlier continuation strengthened necessary conditions for
$\sha{R8_3_3472}$ while its classification was still unresolved.
Theorem~\ref{thm:r3472-general-exclusion} now settles that question.
The statements below retain their original conditional proof domains.
The all-Morita arguments do not assume that a rank-two based module lifts.
The two surviving seven-invertible fusion rings are candidate ring data,
not constructed categories. The parent obstruction below is specific to
the explicit Radchenko--Wheeler family and ordinary de-equivariantization.


% Source: center/center.tex
\subsection{The hypothetical center is simple}
\label{subsec:r3472-focus-simple-center}

Write $R=\sha{R8_3_3472}$ and $r=\sqrt2$. Its positive dimension
character, in canonical order, is
\[
 (1,6+5r,6+5r,7+5r,6+5r,8+5r,8+5r,8+5r),
 \qquad D=700+490r.
\]
An exhaustive check of the $128$ unit-containing subsets shows that
the only based fusion subrings are $\mathbb Z\one$ and $R$.
Independently, every nonunit fusion matrix has entrywise-positive square.
The regular Casimir has characteristic polynomial
\[
 (x-7)^6(x^2-1400x+9800).
\]
Consequently, in a pseudo-unitary categorification, the distinguished
summands of $I(\one)$ have dimensions
\[
 1,\quad 99+70r,\quad
 \underbrace{100+70r,\ldots,100+70r}_{6\text{ copies}}.
\]

\begin{proposition}
If $R$ has a complex categorification $\mathcal C$, then
$Z(\mathcal C)$ is simple. Every nonunit central simple has nonintegral
Frobenius--Perron dimension. In particular its center has no nontrivial
pointed, integral, or Tannakian fusion subcategory.
\end{proposition}
\begin{proof}
Use the previously established general sphericalization and Galois
reduction for $R$. If $\mathcal C$ were braided, its braiding would
transport to its pseudo-unitary Galois conjugate. In a braided category,
the squared norm of a simple is a $d$-number. In the pseudo-unitary
conjugate this is its squared Frobenius--Perron dimension, and an
algebraic integer is a $d$-number if and only if its square is one.
The dimension $6+5r$ has trace $12$ and norm $-14$; the quadratic
$d$-number divisibility condition fails because $14\nmid144$.
Thus no categorification of $R$ is braided.

Liu--Palcoux--Ren \cite[Theorem~5.1]{LPRSimpleFocus} state that a nonpointed simple
complex fusion category has simple center if and only if it admits
no braiding. The theorem does not assume integrality, despite the
title of the paper. The verified based simplicity of $R$ therefore
proves the first assertion.

For the second assertion let $Z$ be a central simple of integral
FP dimension, and write $F(Z)=\sum_i n_iX_i$. The coefficient of $r$
in its FP dimension is $5\sum_{i>0}n_i$. It follows that $F(Z)$ is
a sum of units. Every object in the fusion subcategory generated by
$Z$ consequently also forgets to a sum of units. If $Z\ne\one$,
simplicity would make this the entire center, contrary to the
dominance of the forgetful functor onto the nontrivial category
$\mathcal C$. Hence $Z=\one$.
\end{proof}

These assertions are restrictions on a hypothetical category, not an
exclusion of $R$. In particular the integral prime-power central-simple
obstruction in Remark~5.15 of the cited paper cannot be applied here:
its hypotheses require an integral modular category. The center here
has irrational global FP dimension $D^2$.

The exact ring checks are recorded in
\begin{quote}\small
\path{evidence/r3472_focus_20261022/center/verify_structure.py}.
\end{quote}
They recheck the tensor, based subrings, positive dimension character,
formal codegrees and primitive fusion matrices. They do not replace
the earlier complete signed-quotient proof of the spherical reduction.


\subsection{Exact rational and integer center witnesses}
\label{subsec:r3472-focus-rational-center}

The previously published strengthened center system for
\sha{R8_3_3472} has $3347$ equations and $13242$ variables. Its
constraints comprise the induction Gram matrix, distinguished
$S$-symmetry, cubic Verlinde equations, Galois spectral balance,
Frobenius--Perron bounds, and unit-product constraints. The certificate
\begin{quote}\small
\path{evidence/r3472_focus_20261022/center/rational_witness.json}
\end{quote}
gives an exact nonnegative rational solution satisfying every propagated
box bound. It has $234$ nonzero entries, all integral except one:
the high-column multiplicity for restriction vector
$(0,10,10,10,10,10,10,10)$ equals $49/50$.

The accompanying verifier reconstructs the unchanged published system
from the original census, checks its source inventory and matrix hashes,
replays the box-bound derivation, and verifies all equations and bounds
by rational arithmetic. There is also an exact integer solution to this same system. Let
$u=(0,1,1,1,1,1,1,1)$. The complete column for $10u$ equals $100$ times
the column for $u$. Replacing weight $49/50$ on $10u$ by weight $98$
on $u$ preserves all equations and satisfies all propagated bounds.
Every other coordinate was already integral. The certificate
\begin{quote}\small
\path{evidence/r3472_focus_20261022/center/integer_witness.json}
\end{quote}
records the resulting $234$ nonzero integer entries. The same verifier
checks the full column identity and all equations and bounds for both
witnesses in one reconstruction.

Thus this particular necessary relaxation cannot exclude $R$, even by
integer infeasibility. These certificates do not incorporate twist-phase
equations and do not provide a modular datum or categorification.
At this historical checkpoint categorifiability remained unresolved;
Theorem~\ref{thm:r3472-general-exclusion} supplies the later exclusion.


% Source: morita/manuscript.tex
\subsection{Dimension descent and pointed-orbit restrictions for every Morita partner}
\label{subsec:focus3472-morita}

The following argument does not assume that either rank-two based module
is categorifiable.  It applies to every Morita partner of every complex
categorification of $\sha{R8_3_3472}$.

\begin{lemma}[Dimension-field descent]
\label{lem:focus-dimension-descent}
Let $\mathcal B$ be a complex fusion category such that its Drinfeld
center has no nontrivial invertible objects.  Then
\[
 \Q\bigl(\FPdim(Y):Y\in\Irr(\mathcal B)\bigr)
 =\Q\bigl(\FPdim(\mathcal B)\bigr).
\]
\end{lemma}
\begin{proof}
Write $d_Y=\FPdim(Y)$, $D=\sum_Yd_Y^2$ and $K=\Q(d_Y:Y)$.
Every conjugate of $d_Y^2$ is an eigenvalue of the positive semidefinite
integer matrix $N_YN_{Y^*}=N_YN_Y^{\mathsf T}$; consequently $K$ is
totally real.  Every embedding $\sigma:K\longrightarrow\mathbb R$
satisfies $|\sigma(d_Y)|\leq d_Y$, since $d_Y$ is the spectral radius
of the nonnegative fusion matrix $N_Y$.
If $\sigma(D)=D$, equality in
\[
 D=\sum_Y\sigma(d_Y)^2\leq\sum_Yd_Y^2=D
\]
gives $\sigma(d_Y)=\varepsilon_Yd_Y$, with
$\varepsilon_Y\in\{1,-1\}$. Applying $\sigma$ to the dimension
equation for $XY$, equality in the triangle inequality gives
$\varepsilon_Z=\varepsilon_X\varepsilon_Y$ whenever
$N_{XY}^{Z}>0$. The scalars $\varepsilon_Y$ therefore define a tensor
automorphism of the identity functor, equivalently a central object
whose underlying object is the unit.  The hypothesis forces this
automorphism to be trivial. Thus every embedding of $K$ fixing
$\Q(D)$ is the identity, and separability gives $K=\Q(D)$.
\end{proof}

\begin{proposition}[All-Morita pointed restriction]
\label{prop:focus3472-morita-invertibles}
Let $\mathcal C$ be any complex categorification of
$\sha{R8_3_3472}$ and let $\mathcal B$ be Morita equivalent to
$\mathcal C$. Put $r=\sqrt2$ and $\delta=7+5r$. Then every simple
of $\mathcal B$ has dimension $a+br$, with $a,b\in\NN$, and
\[
 \sum a^2=350,\qquad \sum b^2=175,\qquad\sum ab=245.
 \tag{\ensuremath{\dagger}}\label{eq:focus3472-moments}
\]
The integral-dimensional subcategory is pointed, and its order is
either $1$ or $7$. In the latter case the complete dimension profile
is $1^7,\delta^7$.
\end{proposition}
\begin{proof}
The source tensor has only one invertible basis element, the unit,
and $N_{ii}^{i}=2$ for each of its seven nonunit basis elements.
An invertible object of $Z(\mathcal C)$ must therefore forget to the
unit. Its half-braiding is a tensor automorphism of the identity,
with nonzero scalars $\lambda_i$ satisfying
$\lambda_i^2=\lambda_i$ because $X_i$ occurs in $X_i^2$.
All $\lambda_i$ are one. Thus $Z(\mathcal C)$ has no nontrivial
invertibles, and neither does $Z(\mathcal B)$ by Morita invariance.

Morita invariance also gives
$D=\FPdim(\mathcal B)=700+490r$. The preceding lemma puts every
$d_Y$ in $\Q(r)$, whose ring of integers is $\Z[r]$.
Write $d_Y=a_Y+b_Yr$. Positivity of $d_Y$ and the bound
$|a_Y-b_Yr|\leq a_Y+b_Yr$ imply $a_Y,b_Y\geq0$.
Orthogonality of the two distinct dimension characters gives
$\sum(a_Y^2-2b_Y^2)=0$. Combining this with $\sum d_Y^2=D$
proves \eqref{eq:focus3472-moments}.

The integral-dimensional simples form a fusion subcategory because
all radical coefficients are nonnegative. Let $h$ be its dimension.
The integral-dimension Galois Gram bound gives $h\leq7$, since
$\operatorname{Tr}_{\Q(r)/\Q}(D^{-1})=1/7$.
Categorical Lagrange integrality of $D/h$ gives
$1400/h,9800/h^2\in\Z$, whence $h\in\{1,2,5,7\}$.
This subcategory is pointed. The only nonpointed integral profiles
at these dimensions are $1,2$ at dimension five and $1,1,1,2$ at
dimension seven. Their unique dimension-two simple is self-dual and
fixed by every invertible, giving respectively the impossible
dimension equations $4=1+2m$ and $4=3+2m$.
Write $g=h$ for the number of invertibles.

We use strong Lagrange divisibility: if $\mathcal D\subseteq\mathcal B$
is a fusion subcategory and $\mathcal M$ is an indecomposable
$\mathcal D$-module component of $\mathcal B$, then
$\FPdim(\mathcal M)/\FPdim(\mathcal D)$ is an algebraic integer
\cite[Proposition~5.7]{IntegralPublished}. For the pointed subcategory and a
left orbit $\mathcal O$ of simples of common dimension $d$, this says
$|\mathcal O|d^2/g$ is integral.
If $g=2$, a fixed simple of dimension $a+br$ must have
$d^2/2$ integral, so $a$ is even. Free orbits contribute an even
integer to $\sum ab$, as do fixed simples. This contradicts
$\sum ab=245$.

Suppose $g=5$. A fixed simple has $d^2/5$ integral, so
$a^2+2b^2\equiv 2ab\equiv0\pmod5$, implying $5\mid a,b$.
From \eqref{eq:focus3472-moments}, after removing the invertibles,
\[
 \sum(5a-7b)^2=50.
\]
Together with $a^2\leq345$, $b^2\leq175$ this leaves only
$(a,b)=(15,10)$ for a fixed nonintegral simple. There can be at most
one. After removing it, the remaining nonintegral dimensions would
satisfy $\sum a^2=120$, $\sum b^2=75$, $\sum ab=95$, contrary to
$95^2>120\cdot75$. Hence every nonintegral orbit has length five.
For orbit representatives, the moment equations become
\[
 \sum a^2=69,\qquad\sum b^2=35,\qquad\sum ab=49,
 \qquad\sum(5a-7b)^2=10.
\]
The complete possible types are
$(1,1),(2,1),(3,2),(4,3),(5,4),(6,4),(7,5)$.
Their nonnegative multiplicities give exactly the following four
lists of nonintegral orbit-representative dimensions:
\[
\begin{array}{c|l}
 \mathrm A&2+r,\ 4+3r,\ 7+5r\\
 \mathrm B&1+r,\ (4+3r)^2,\ 6+4r\\
 \mathrm C&1+r,\ (3+2r)^4,\ (4+3r)^2\\
 \mathrm D&(1+r)^2,\ (3+2r)^2,\ 7+5r.
\end{array}
\]
Superscripts denote multiplicities of free orbits. This finite
seven-type enumeration is replayed with exact integer arithmetic.
In A, the product $XX^*$ for $\FPdim(X)=2+r$ contains exactly one
invertible, since its stabilizer is trivial. Its nonintegral part
would have dimension $5+4r$, which is no nonnegative combination of
the three available dimensions (their rational parts are even or
at least seven). Thus A is impossible.

For B and C, put $u=1+r$. There is one free orbit of $u$-dimensional
simples, necessarily stable under duality. The right orbit of $X^*$
is free, since its stabilizer is dual to the left stabilizer of $X$,
and therefore fills this same five-element orbit. Hence for any two
$u$-dimensional simples $X,Y$ there is exactly one $g\in C_5$ with
$Y\simeq X^*g$. Their product contains exactly one invertible
and has nonintegral part $u^2-1=2u$. No larger available dimension
fits that remainder. They therefore generate, with the five
invertibles, a fusion subcategory of profile $1^5,u^5$.

In D, there are two $u$-orbits, two $u^2$-orbits and one $u^3$-orbit.
Right multiplication by $C_5$ preserves each left orbit: it acts on
the left orbits of any fixed dimension, and there are fewer than
five of these. Hence identifying the five invertibles with the
unit gives a based fusion-ring quotient on the six left orbits,
of dimensions $1,u,u,u^2,u^2,u^3$. Here are the details that ensure
this is a based quotient, rather than merely a dimension quotient.
For orbit representatives $x,y,z$, set
$\overline N_{[x][y]}^{[z]}=\sum_{g\in C_5}N_{xy}^{gz}$.
Left translation of $x$ translates the product, and replacing $y$
by $gy$ does the same because $xg$ belongs to the left orbit of $x$.
The coefficients are therefore independent of representatives and
inherit associativity and duality. Their unit coefficient counts
$g$ with $y\simeq x^*g$, so it is one precisely for dual orbits.
Finally, Frobenius reciprocity follows from
\[
 \sum_gN_{xy}^{gz}
 =\sum_gN_{y,z^*g^{-1}}^{x^*}
 =\sum_gN_{y,z^*}^{x^*g}
 =\overline N_{[y][z^*]}^{[x^*]},
\]
since right multiplication runs through the same free left orbit.
Write $A,B$ for the two $u$-dimensional basis elements of this quotient.
A product with a unit term has the form $1+$ two $u$-terms;
without a unit term it is a single $u^2$-dimensional basis element.
If $A^*=B$, the squares $A^2,B^2$ have no $u$-terms. Cyclic
Frobenius reciprocity then gives
$[AB:A]=[BB:B]=0$ and $[AB:B]=[AA:A]=0$, contradicting
$\FPdim(AB-1)=2u$. Thus both are self-dual. Now $AB$ is one
$u^2$-dimensional basis element, so reciprocity gives
$[A^2:B]=[AB:A]=0$ and similarly $[B^2:A]=0$. Consequently
$A^2=1+2A$ and $B^2=1+2B$. Lifting either orbit back to
$\mathcal B$ again gives a fusion subcategory of profile $1^5,u^5$.
Lemma~\ref{lem:focus-c5-silver} below excludes this profile and
finishes the order-five case.

Finally, if $g=7$, equality in the Galois Gram bound forces
$5a=7b$ for every nonintegral dimension. Thus those dimensions are
$k\delta$, with positive integers $k$ satisfying $\sum k^2=7$.
The possibilities are $\delta^7$ and $\delta^3,(2\delta)^1$.
Since $\delta^2=1+14\delta$, a $\delta$-dimensional simple has
trivial stabilizer among the seven invertibles. Its orbit has seven
members, excluding the second possibility.
\end{proof}

\begin{lemma}[The order-five silver-ratio profile]
\label{lem:focus-c5-silver}
There is no complex fusion category with dimension profile
$1^5,(1+\sqrt2)^5$ whose invertibles act freely on the
noninvertible simples.
\end{lemma}
\begin{proof}
Let $g$ generate the invertibles and choose a noninvertible $X$.
The left action is free and transitive. The right action is too:
the coefficients of invertibles in $XX^*$ and $X^*X$ are each
forced by $\FPdim(X)^2=1+2\FPdim(X)$ to sum to one. Thus
$Xg=g^aX$ for some $a\in(\Z/5\Z)^\times$.
Since the dual remains in the same free orbit, dualizing the
relation gives $a^2=1$, so $a=\pm1$.
One can choose $X=X^*$: for $a=1$ shift its index by half its
duality translation, and for $a=-1$ the involution forces that
translation to vanish. Write
\[
 X^2=1+\sum_{j\in\Z/5\Z}m_jg^jX,
 \qquad m_j\in\NN,\qquad\sum m_j=2.
\]
If $a=-1$, comparing $X^2g$ with $X(Xg)$ forces $m_j=m_{j+2}$,
which would make all five $m_j$ equal, an impossibility.
If $a=1$, duality gives $m_j=m_{-j}$. There are two possibilities
up to changing the generator: $m_0=2$, or $m_1=m_{-1}=1$.
The first contains the rank-two fusion subcategory
$X^2=1+2X$, excluded by the rank-two classification
\cite{Ostrik2}.

In the second case the ring is commutative. Choose a primitive fifth
root $\zeta$ and a one-dimensional character with
$\chi(g)=\zeta$, $\chi(X)=x$, where
$x^2-(\zeta+\zeta^{-1})x-1=0$. Its formal codegree is
$f=5(1+x^2)$. Since $t=\zeta+\zeta^{-1}$ satisfies
$t^2+t-1=0$, eliminating $x,t$ gives
\[
 f^4-55f^3+1000f^2-7250f+18125=0.
\]
The polynomial is irreducible: modulo two it is $z^4+z^3+1$,
which has neither a root nor the factor $z^2+z+1$.
The formal-codegree $d$-number criterion
\cite[Theorem~1.2 and Lemma~2.7]{OstrikCode} would
require $18125\mid55^4$, whereas the remainder is $15625$.
This excludes the second case without a pivotality assumption.
\end{proof}

\begin{corollary}
The first rank-two branch of
Proposition~\ref{prop:creative-r3472-three} retains only four of its
seven dimension spectra: its three one-invertible spectra and the
seven-invertible spectrum $1^7,(7+5\sqrt2)^7$.
\end{corollary}
\begin{proof}
The all-Morita proposition removes the remaining order-two spectrum
and both order-five spectra. This is still a conditional restriction
on the rank-two construction route: neither its lift nor a
categorification of the source ring is asserted.
\end{proof}

The verifier in
\path{evidence/r3472_focus_20261022/morita/verify.py} checks the
source tensor, all arithmetic bounds, the complete four-profile
enumeration and the irreducible formal-codegree obstruction.
The general categorification status of $\sha{R8_3_3472}$ remains open.


% Source: constructions/constructions.tex
\subsection{The seven-invertible Morita candidates}
\label{subsec:focus-r3472-candidates}

Put $r=\sqrt2$, $\delta=7+5r$, and $D=700+490r$.
The following statement concerns the seven-invertible alternative in the
preceding Morita restriction. It does not assert that this alternative occurs.

\begin{proposition}\label{prop:focus-r3472-two-rings}
Every seven-invertible Morita partner of a categorification of
$\sha{R8_3_3472}$ has one of the following two based rings. Their basis is
$\{g_i,X_i:i\in\mathbb Z/7\mathbb Z\}$, with
\[
 g_i g_j=g_{i+j},\qquad g_iX_j=X_{i+j},\qquad
 X_i g_j=X_{i+aj},\qquad
 X_iX_j=g_{i+aj}+2\sum_kX_k,
 \quad a\in\{1,-1\}.
\]
Both displayed multiplications are associative fusion rings. No associator
or Morita identification with the source ring is established here.
\end{proposition}
\begin{proof}
The preceding dimension restriction gives $1^7,\delta^7$.
Since $\delta^2=1+14\delta$, reciprocity makes both stabilizers of every
noninvertible trivial. Write $X_i=g^iX$ and $Xg=g^aX$, where $a$ is a
unit modulo seven. The unique invertible summand of $X^2$ is $g^b$.
Comparing the invertible parts in $X^2g=g^{a^2}X^2$ gives $a^2=1$.
Reciprocity gives $X^*=g^{-b}X$, and applying duality twice gives
$(a-1)b=0$. Thus $b=0$ if $a=-1$, and replacing $X$ by $g^cX$ makes
$b=0$ if $a=1$. In particular $X$ is self-dual.

For $a=-1$, the identity $X^2g=gX^2$ makes the noninvertible
coefficients invariant under translation by two. They sum to fourteen,
so all are two. For $a=1$, write
$X^2=1+m(g)X$ with $m_j=m_{-j}\geq0$ and $\sum_jm_j=14$.
Pass to the pseudo-unitary Galois conjugate supplied by the source
analysis; this does not change any fusion coefficient. All simple
dimensions in its common center belong to $\mathbb Q(r)$, by forgetting
to the source category. The formal-codegree induction formula
$\operatorname{FPdim}(A_E)=D/f_E$ therefore gives
$f_E\in\mathbb Q(r)$ for every irreducible representation $E$.

Choose a primitive seventh root $\zeta$. The characters above
$g\mapsto\zeta$ have $X\mapsto x$, where
$x^2-tx-1=0$ and $t=m(\zeta)$ is real. Their formal codegrees are
$7(1+x^2)$. Consequently
\[
 t^2=\frac{(x^2-1)^2}{x^2}
 \in\mathbb Q(r)\cap\mathbb Q(\zeta+\zeta^{-1})=\mathbb Q.
\]
The latter field has odd degree three, hence $t\in\mathbb Q$.
The polynomial $\Phi_7$ now forces $m_0=A$ and $m_j=B$ for $j\ne0$,
with $A+6B=14$. The possibilities are $(A,B)=(14,0),(8,1),(2,2)$.
The first contains the impossible rank-two categorification
$X^2=1+14X$. The second has formal codegrees
$(371\pm49\sqrt{53})/2$, outside $\mathbb Q(r)$. Hence $A=B=2$.
The complete tensors, dualities, dimensions, reciprocity identities and
all $14^4$ associativity equations for each ring are checked in
\path{evidence/r3472_focus_20261022/constructions/verify_constructions.py}.
\end{proof}

\subsection{An obstruction to all direct Radchenko--Wheeler parent routes}
\label{subsec:focus-r3472-rw}

We use the explicit near-group categories of Radchenko--Wheeler,
\emph{Real quadratic fields and finite quantum dilogarithms I},
\cite[Sections~4.3 and~5.4]{RWQuadraticFocus}. The scope below is their
construction from a matrix $\gamma\in\mathrm{SL}_2(\mathbb Z)$;
no classification of all categorifications of the parent fusion ring
is asserted.

\begin{lemma}\label{lem:focus-rw-halfbraiding}
For a Radchenko--Wheeler matrix input with
$\operatorname{tr}(\gamma)=N+2$, put
\[
 G=\mathbb Z^2/\mathbb Z^2(\gamma-I),\qquad
 \chi(u,v)=\exp\!\left(\frac{2\pi i}{N}
                  \omega(u(\gamma-I),v)\right).
\]
The central invertibles embed into $G$, and their self-braidings agree
with $q(g)=\chi(g,g)^{-1}$ on their image.
\end{lemma}
\begin{proof}
The strict algebraic construction gives
$\alpha_h\rho=\rho$,
$\rho\alpha_g=\operatorname{Ad}(U_g)\rho$, and
$\alpha_h(U_g)=\chi(h,g)U_g$.
The half-braiding of a central lift of $\alpha_g$ with $\rho$ is a
nonzero scalar times $U_g$, since $\rho$ is simple. Its hexagon at
$\alpha_h\rho=\rho$ forces its half-braiding with $\alpha_h$ to be
$\chi(h,g)^{-1}$. This uses neither a star operation nor unitarity.
The forgetful kernel on central invertibles consists of tensor
automorphisms of the identity. It is trivial because the universal
grading is trivial: $\rho$ and every $\alpha_g$ occur in $\rho^2$.
\end{proof}

\begin{lemma}\label{lem:focus-rw-odd-square}
For $N=4s^2$ with $s$ odd, the bosonic part of the radical of
$(G,q)$ has order at most two.
\end{lemma}
\begin{proof}
The polarization is $\chi^{-2}$, with radical $G[2]$.
The two-primary part of the two-generated group $G$ has order four.
If it is cyclic, $|G[2]|=2$. Otherwise
$\gamma-I=2A$, where
\[
 \operatorname{tr}A=2s^2\equiv2\pmod4,\qquad
 \det A=-s^2\equiv3\pmod4.
\]
The two rows of $A$ represent a basis of $G[2]$. For
$A=\left(\begin{smallmatrix}a&b\\c&d\end{smallmatrix}\right)$,
the corresponding diagonal bicharacter values are $(-1)^b$ and
$(-1)^c$. Triviality of $q|_{G[2]}$ would make $b,c$ even.
Then $a,d$ are odd and their sum is two modulo four, so $ad$ is one
modulo four. This contradicts $\det A\equiv3\pmod4$.
Thus $q$ is a nontrivial character on the four-element radical,
and its kernel has order two.
\end{proof}

\begin{theorem}\label{thm:focus-rw-pointed-stabilization}
Let $s$ be odd and positive, set
$\delta_s=s+\sqrt{s^2+1}$ and $D_s=s(1+\delta_s^2)$.
No Radchenko--Wheeler parent with group order $4s^2$, tensored with
an arbitrary pointed fusion category, admits an ordinary
de-equivariantization of dimension $D_s$.
In particular, no such direct route, including noncyclic parents and
pointed tensor factors, supplies a Morita representative for
$\sha{R8_3_3472}$.
\end{theorem}
\begin{proof}
Let the pointed factor $\mathcal P$ have dimension $m$.
The parent has dimension $4sD_s$, so the required Tannakian group
would have order $h=4sm$. Its fully faithful forgetful image is pointed:
simple representations have integer dimensions, whereas all
noninvertible dimensions in the parent tensor product are irrational.
Thus its simples form an isotropic subgroup $H$ of order $h$ in the
central invertibles.

Put $L=\operatorname{Inv}(Z(\mathcal P))$, $k=|L|$, and let the
radical of its self-braiding form have order $t$.
The centralizer dimension formula in the nondegenerate category
$Z(\mathcal P)$ gives $kt\leq m^2$.
Lemma~\ref{lem:focus-rw-halfbraiding} embeds $H$, preserving its
quadratic form, into $G\times L$. Its intersection with the radical
has order at most $2t$: in the cyclic two-primary case the full
radical has that order, and in the other case the radical quadratic
form is nontrivial already on $G[2]$, so its kernel has that order.
For any isotropic subgroup of a finite premetric group $A$,
\[
 |H|^2\leq |A|\,|H\cap\operatorname{rad}(A)|,
\]
because $H\subseteq H^\perp$ and
$|H^\perp|=|A|\,|H\cap\operatorname{rad}(A)|/|H|$.
It follows that
$h^2\leq4s^2k(2t)\leq8s^2m^2$, contradicting
$h^2=16s^2m^2$.

For the stated target, no other near-group parent order can have
the required rational dimension scale. The global dimension $E_N$
of $\mathrm{NG}(G,N)$, $|G|=N$, has trace $N(N+4)$ and norm
$N^2(N+4)$. Hence its scale-invariant trace-squared/norm ratio is
$N+4$. For $D$ the ratio is $1400^2/9800=200$, forcing $N=196$.
The case $s=7$ above therefore covers all Radchenko--Wheeler matrix
parents with a possible dimension match.
\end{proof}

This rules out the specified direct ordinary quotients, also followed by
a Morita equivalence. It does not address a prior Morita change of the
parent, arbitrary near-group categorifications with different input data,
generalized-Haagerup parents, or arbitrary chains of constructions.
The two rank-fourteen tensors remain unconstructed candidate based rings;
$\sha{R8_3_3472}$ was unresolved at this checkpoint. Its subsequent
general exclusion is Theorem~\ref{thm:r3472-general-exclusion}.


% Source: independent/third_indicators.tex
\subsection{Maximal third indicators}
\begin{proposition}\label{prop:r3472-maximal-third-indicator}
In the canonical positive spherical structure of every pseudo-unitary
categorification of $\sha{R8_3_3472}$,
the third Frobenius--Schur indicator of every nonunit simple is $2$.
Equivalently, rotation acts identically on each two-dimensional space
$\operatorname{Hom}(1,X_i^{\otimes3})$, $i>0$.
\end{proposition}
\begin{proof}
Put $\delta=7+5\sqrt2$, $D=700+490\sqrt2$, and $D'=700-490\sqrt2$.
The source dimensions are
$(1,\delta-1,\delta-1,\delta,\delta-1,\delta+1,\delta+1,\delta+1)$,
and $\dim\operatorname{Hom}(1,X_i^{\otimes3})=2$ for all $i>0$.
The spherical Cauchy theorem~\cite[Theorem~3.9]{BNRWCauchy} and
$DD'=9800$ imply that $3$ is coprime to
$N=\operatorname{FSexp}(\mathcal C)$.
The dimensions lie in $\mathbb Q(\zeta_N)$, so $8\mid N$.
For $3\widetilde3\equiv1\pmod N$, the automorphism
$\sigma(\zeta_N)=\zeta_N^{\widetilde3}$ therefore flips $\sqrt2$.
The higher Gauss-sum formula~\cite[Theorem~4.4, (4.17)]{NSWHigherGauss}
gives
\[
 \sum_i d_i\nu_3(X_i)=\frac D{\sigma(D)}
 =\frac D{D'}=\delta^2=1+2\sum_{i>0}d_i.
\]
We have $\nu_3(1)=1$. For $i>0$, $\nu_3(X_i)$ is the trace of an
order-three operator on a two-dimensional space, so
$\operatorname{Re}\nu_3(X_i)\le2$.
All $d_i$ are positive, so the displayed equality forces equality
for every summand. Every rotation eigenvalue is consequently $1$;
the operator is diagonalizable because its cube is the identity.
\end{proof}

The existing height-four full-character-field phase witness already
has third-rotation weights $(0,2)$, for trace types $(-1/2,1)$,
at every nonunit simple. Thus this restriction does not exclude that
relaxation. Exact source arithmetic and the cached-coordinate audit are
recorded in
\path{evidence/r3472_focus_20261022/independent/verify_third_indicators.py}.


\paragraph{An exact phase extension of the integral center point.}
The integral center-relaxation point also admits a nonnegative rational
extension satisfying all relative indicator and rotation moments through
degree four, with the third indicators fixed by
Proposition~\ref{prop:r3472-maximal-third-indicator}.
Its ten positive high restriction types have multiplicities
$2,\ldots,2,98$, the last belonging to $(0,1,1,1,1,1,1,1)$.
Splitting these counts among the complete $37$ joint relative-trace
types, and adding three fourth-rotation variables for each nonunit
simple, gives $391$ variables. The $80$ equations consist of ten
fixed-mass equations, $42$ coordinate equations for indicators of
degrees $1,2,3$, fourteen fourth-rotation dimension/square-trace
equations, and fourteen fourth-indicator coordinate equations.
An exact rational point has $40$ positive entries.

The source-bound verifier
\begin{quote}\small\path{evidence/r3472_focus_20261022/independent/verify_phase_extension.py}\end{quote}
regenerates the trace domain and checks every equation using rational
arithmetic. It binds the restriction labels to the independently
verified integral center point, and passes both normally and under
\texttt{python -O}. The trace compression bound is $1120$; orders
outside its divisor box have zero joint traces through degree four,
so no bound on the categorical FS exponent is assumed.
The added constraints are indicator/rotation moments, not the complete
complex modular relation $STS$. The phase and rotation weights remain
rational; this feasible necessary relaxation supplies neither actual
modular data nor a categorification.

\subsection{Recovery of the arithmetic-count verifier}
The preceding source package omitted the imported counting helper
\path{second_branch_count.py}. Its reconstructed implementation reproduces
exactly $272348$, $124228$, $3900$, and $2$ second-branch spectra for
$h=1,2,5,7$, respectively. A separate generating-function algorithm
independently verifies these counts. The three original mathematical
scripts are unchanged. The historical receipt is preserved, and the new
receipt differs only in four cache-state diagnostics. Both normal and
optimized executions pass; no mathematical assertion was weakened.
The repair, original source bindings, and independent replay are in
\path{evidence/r3472_focus_20261022/reproducibility/}.
These historical arithmetic spectra are superseded, as necessary Morita
data, by the all-Morita restrictions above; none represents an asserted
categorical construction.

For the second branch, the new order-two and order-five exclusions remove
$124228$ and $3900$ of the previously retained $400477$ arithmetic spectra.
Exactly $272349=272348+1$ dimension multisets remain under these filters:
the one-invertible spectra and the sole seven-invertible profile
$1^7,\delta^7$. Indeed $\delta=7+5r$ has the old coordinates
$(s,t,u)=(7,5,2)$, whose seven copies give
$\sum u^2=28$, $\sum ut=70$, and $\sum t^2=175$.
This is a count of necessary arithmetic profiles, not of fusion rings
or fusion categories.
